\documentclass[10pt,letterpaper]{article}
\usepackage[margin=0.65in]{geometry}
\usepackage[T1]{fontenc}
\usepackage{amsmath,amssymb,enumitem,xcolor,hyperref,microtype,titlesec}
\usepackage[backend=biber,style=authoryear,maxcitenames=1,maxbibnames=1,giveninits=true,uniquename=false,doi=false,url=false]{biblatex}
\addbibresource{references.bib}
\hypersetup{colorlinks=true,citecolor=black,urlcolor=blue!55!black}
\setlength{\parindent}{0pt}
\setlength{\parskip}{3pt}
\titleformat{\section}{\normalsize\bfseries}{}{0pt}{}
\titlespacing*{\section}{0pt}{8pt}{3pt}
\setlist[itemize]{leftmargin=*,nosep,topsep=2pt}
\renewcommand*{\bibfont}{\footnotesize}
\setlength{\bibitemsep}{1pt}
\begin{document}

We reproduced the experimental setup by Arditi et. al for evaluation, with the following differences:
\begin{itemize}
    \item Although we use the same datasets, we sample from them differently to make sure that it is possible to reproduce their results with a different random sample.
    \item We replace their direction selection algorithm with a simpler approach that chooses the refusal direction from the layer that most effectively separates refusal from non-refusal using the decision boundary that is equidistant from the mean harmful activation and mean harmless activation.
    \item 
\end{itemize}
\end{document}
